\documentclass[aspectratio=169]{beamer}

% Tema y configuración
\usetheme{Madrid}
\usecolortheme{default}

% Paquetes necesarios
\usepackage[utf8]{inputenc}
\usepackage[spanish]{babel}
\usepackage{amsmath}
\usepackage{amsfonts}
\usepackage{amssymb}
\usepackage{graphicx}
\usepackage{listings}
% Habilitar UTF-8 dentro de lstlisting (acentos, ñ, etc.)
\usepackage{listingsutf8}
\usepackage{xcolor}
\usepackage{verbatim}
% Para dibujos con TikZ
\usepackage{tikz}
\usetikzlibrary{positioning,trees}

% Configuración de listings para código
\lstset{
    basicstyle=\ttfamily\footnotesize,
    keywordstyle=\color{blue}\bfseries,
    commentstyle=\color{green},
    stringstyle=\color{red},
    showstringspaces=false,
    breaklines=true,
    frame=single,
    numbers=left,
    numberstyle=\tiny\color{gray},
    inputencoding=utf8,
    % Mapeo de caracteres españoles para que no falle la compilación
    literate={á}{{\'a}}1 {é}{{\'e}}1 {í}{{\'\i}}1 {ó}{{\'o}}1 {ú}{{\'u}}1
             {Á}{{\'A}}1 {É}{{\'E}}1 {Í}{{\'I}}1 {Ó}{{\'O}}1 {Ú}{{\'U}}1
             {ñ}{{\~n}}1 {Ñ}{{\~N}}1 {ü}{{\"u}}1 {Ü}{{\"U}}1
}

% Información del documento
\title{Hoja 5: Ejercicio 23}
\subtitle{TEMA}
\author{Diego Rodríguez Cubero}
\institute{UCM}
\date{\today}

% Logo (opcional)
% \logo{\includegraphics[height=1cm]{images/logo.png}}

\begin{document}

% Diapositiva de título
\begin{frame}
    \titlepage
\end{frame}

% Tabla de contenidos
\begin{frame}{Contenidos}
    \tableofcontents
\end{frame}

% Sección 1
\section{Enunciado del problema}
\begin{frame}{Enunciado del problema}
    \begin{block}{Ejercicio 23}
    El dueño de una pinacoteca quiere comprar nuevos cuadros para incorporarlos a su colección, pero no desea comprar más de un cuadro por mes. El precio de los n cuadros que le interesan en el momento actual es conocido $(p_1, \ldots , p_n)$ pero a medida que pasan los meses los precios se irán incrementando de manera exponencial, porque se trata de un mercado con una gran competencia. Así que, si el precio inicial es $p_i$ , después de $t$ meses el precio será $p_i^t$. Empezará comprando los cuadros al final del primer mes $(t = 1)$ y sucesivamente al final del segundo mes, del tercero, etc. Por ejemplo, si $n = 3$, los cuadros cuestan $p_1 = 3, p_2 = 5$ y $p_3 = 2$ euros, y decide comprar el cuadro i al final del mes i, pagará en total $31 + 52 + 23 = 36$ euros. Explicar e implementar un algoritmo voraz para minimizar el coste total de los n cuadros que quiere comprar. Justificar la corrección de dicho algoritmo. Analizar el coste asintótico del algoritmo implementado.
    \end{block}
\end{frame}

% Sección 2
\section{Algoritmo voraz propuesto}
\begin{frame}{Algoritmo voraz propuesto}
    \begin{enumerate}
        \item Ordenar los precios de los cuadros en orden descendente.
        \item Comprar los cuadros en ese orden, asignando el mes correspondiente.
    \end{enumerate}
\end{frame}

\begin{frame}[fragile]{Implementación en C++}
    \begin{lstlisting}[language=C++,title=Algoritmo]
    int coste_Minimo(vector<int> &precios) {
    sort(precios.begin(), precios.end(), greater<int>()); // Ordenamos de mayor a menor
    int n = precios.size();
    int coste_total = 0;
    for (int i = 0; i < n; ++i) {
        coste_total += pow(precios[i], i + 1); // Compramos el cuadro i en el mes i+1
    }
    return coste_total; // Devolvemos el coste total
    }
    \end{lstlisting}
\end{frame}

% Sección 3
\section{Justificación de la corrección}
\begin{frame}{Justificación de la corrección}
    Comprobemos que el algoritmo es correcto mediante un argumento de intercambio, llamemos $X_v$ a la solución voraz y $X_o$ a una solución óptima diferente, siendo ambos casos un vector con el orden de compra de los cuadros, y del mismo tamaño $n$ (número de cuadros).
    \begin{enumerate}
        \item Si $X_v = X_o$, entonces la solución voraz es óptima.
        \item Si $X_v \neq X_o$, entonces tomamos $k = \min \{ i | X_v[i] \neq X_o[i] \}$, es decir, el primer índice donde difieren ambas soluciones (existe al menos uno ya que son diferentes, y todos los cuadros se deben comprar).
        \item Intercambiamos en $X_o$ el cuadro que se compra en la posición $k$ por el cuadro que se compra en la posición $j > k$ donde se encuentra el cuadro que en $X_v$ se compra en la posición $k$.
        \item Llamemos $X_o'$ a la nueva solución tras el intercambio. Entonces, el coste de $X_o'$ es menor o igual que el coste de $X_o$, ya que hemos adelantado la compra de un cuadro más caro y retrasado la compra de un cuadro más barato.
    \end{enumerate}
\end{frame}

\begin{frame}{Justificación de la corrección (cont.)}
    \begin{enumerate}
        \setcounter{enumi}{4}
        \item Repetimos el proceso hasta que $X_o$ sea igual a $X_v$. En cada paso, el coste de la solución no aumenta, por lo que el coste de $X_v$ es menor o igual que el coste de $X_o$.
        \item Por lo tanto, la solución voraz $X_v$ es óptima.
    \end{enumerate}
\end{frame}

\begin{frame}{Justificación de la corrección (cont.)}
    El paso $4$ es válido ya que el coste de un cuadro $i$ comprado en el mes $t$ es $p_i^t$, y si adelantamos la compra de un cuadro más caro y retrasamos la compra de un cuadro más barato (el cuadro que se compra en la posición $k$ de $X_o$ es mas caro que el que se compra en la posición $j > k$, por la ordenación que usa $X_v$ y ya que el cuadro $X_o[k]$ todavia no es comprao en $X_v$), el coste total disminuye o se mantiene igual, usando formulas:
    \[
    p_{X_o[k]}^k + p_{X_o[j]}^j \leq p_{X_o[k]}^j + p_{X_o[j]}^k \iff
    p_{X_o[k]}^k - p_{X_o[j]}^k \leq p_{X_o[k]}^j - p_{X_o[j]}^j \iff
    \]
    \[
    (p_{X_o[k]} - p_{X_o[j]}) \sum_{i=0}^{k-1} p_{X_o[k]}^i p_{X_o[j]}^{k-1-i} \leq (p_{X_o[k]} - p_{X_o[j]}) \sum_{i=0}^{j-1} p_{X_o[k]}^i p_{X_o[j]}^{j-1-i}
    \]
    Que se cumple ya que $p_{X_o[k]} > p_{X_o[j]}$ y la suma de la derecha es mayor que la de la izquierda, al poder asociar a cada termino de la izquierda un término mayor en la derecha:
    \[
    p_{X_o[k]}^{i} p_{X_o[j]}^{k-1-i} \leq p_{X_o[k]}^{i} p_{X_o[j]}^{j-1-i} \quad \forall i \in [0, k-1]
    \]
\end{frame}

% Sección 4
\section{Análisis del coste asintótico}
\begin{frame}{Análisis del coste asintótico}
    \begin{itemize}
        \item El coste de ordenar los precios es $O(n \log n)$, donde $n$ es el número de cuadros.
        \item El coste de calcular el coste total es $O(n)$, ya que recorremos la lista de precios una vez.
        \item Por lo tanto, el coste total del algoritmo es:
    \end{itemize}
    \[
    O(n \log n) + O(n) = O(n \log n)
    \]
    Así, el algoritmo voraz propuesto tiene un coste asintótico de $O(n \log n)$.
\end{frame}

% Sección 5
\section{Conclusión}
\begin{frame}{Conclusión}
    Hemos propuesto un algoritmo voraz para minimizar el coste total de la compra de cuadros en una pinacoteca, basado en ordenar los precios de los cuadros en orden descendente y comprarlos en ese orden. Hemos justificado la corrección del algoritmo mediante un argumento de intercambio y hemos analizado su coste asintótico, concluyendo que es $O(n \log n)$.
\end{frame}

\end{document}